\section{Execution Architecture and Incremental Modularization}
\label{sec:execution-architecture}

GAWorld combines a connected cognition pipeline with a small kernel, nine
built-in plugins, and compatibility-preserving domain paths.
Figure~\ref{fig:system-architecture} freezes the audited K1--K5 runtime at
repository commit \texttt{35a551e}.  The archived capability cases predate
this migration and are not evaluations of it.

\begin{figure}[t]
  \centering
  \includegraphics[width=\linewidth]{system_architecture.pdf}
  \caption{Audited execution architecture at commit \texttt{35a551e}.  Nine
  built-ins contribute through the registry, bus, pipeline filters, or
  lifecycle events.  Movement uses the Controller validation gate; Recorder
  events coexist with established domain-specific outputs.}
  \label{fig:system-architecture}
\end{figure}

\subsection{Kernel services}

\paragraph{Simulation context and clock.}
\texttt{SimContext} holds configuration, Clock, EventBus, PluginRegistry,
Controller, Recorder, agents and their identifier index, an LLM callable, a
seeded random-number generator, and an \texttt{extras} compatibility escape
hatch.  It also provides simulation-level and per-agent plugin-state
namespaces.  The outer loop is the sole clock writer, giving stages, plugins,
and records a common simulated day, time, and tick.  This ordering does not by
itself make model calls, external services, or concurrency deterministic.

\paragraph{Event bus and plugin registry.}
The EventBus supports priority-ordered observation (\texttt{emit}), merged
contribution (\texttt{collect}), and value rewriting (\texttt{filter}).
Handler failures are logged and isolated by default, so warnings and active
plugin lists must accompany results.  The PluginRegistry assembles built-ins,
configuration-named classes, and installed entry points; orders setup by
declared dependencies; skips failed or unsatisfied plugins; and tears active
plugins down in reverse order.  Plugins may own configuration, namespaced
agent state, and output directories, but execute in process and are not a
sandbox.

\paragraph{Controller.}
The Controller provides a priority-ordered validator chain over structured
\texttt{ActionRequest}s and a named runtime-intervention API.  In the main
simulator, movement requests pass through validation: denial preserves the
origin and is surfaced in the next perception, while a rewrite can change the
destination.  K5 registers standard interventions for setting one agent-state
value, updating a dotted live-config path, and queueing agent removal for the
next day boundary.  Denials and interventions are audited through the
Recorder.  Other selected actions are not all mediated, so the Controller is
not a universal policy or safety gate.

\paragraph{Recorder.}
The Recorder appends clock-stamped JSON objects to table-specific files and is
closed when the simulation exits.  It records Controller activity and is
available to plugins.  Main traces, state history, memories, environment
events, and specialized metrics also retain established writers.  Recorder
integration therefore improves cross-plugin auditability without implying
that the archive can be replayed from one event log.

\subsection{Built-in plugin assembly}

The audited registry assembles Intervention, Skills, Interests, Life Events,
Economy, Local Physical, Real Work, Dynamic Behavior, and Spatial Preferences.
They attach domain behavior to day/tick events, perception contributions,
location filters, action or outcome filters, and post-step or day-end hooks.
For example, Intervention contributes recommendation context and persists
policy metrics, Life Events produces scheduled exogenous events and state
effects, Local Physical contributes nearby conditions, and Spatial Preferences
can rewrite a resolved destination.  The adapters delegate calculations and
specialized persistence to established modules: migration changes
orchestration and ownership, not the scientific validity of the mechanisms.

\subsection{Migration-state audit}

\begin{table}[t]
\centering
\small
\setlength{\tabcolsep}{4pt}
\begin{tabular}{@{}p{0.17\linewidth}p{0.23\linewidth}p{0.19\linewidth}p{0.32\linewidth}@{}}
\toprule
Surface & Current boundary & Status & Retained boundary and consequence \\
\midrule
Cognition & Configured 12-stage pipeline & Wired & Shared step dictionary; order and omissions must be reported \\
Extensions & EventBus and PluginRegistry & Wired & Failures may degrade rather than terminate a run \\
Nine domain adapters & Policy, skills, interests, events, economy, physical, work, behavior, spatial & Plugin-integrated & Established stores and writers remain; activation needs config/plugin provenance \\
Memory & Recall, storage, consolidation, decay & Direct compatibility path & Persistent context is inspectable, not psychologically validated \\
Environment and social network & World evolution and network modules & Direct compatibility path & The full world and network models are not isolated plugins \\
Controller & Move validation and named interventions & Connected, scoped & Other actions retain established paths; no universal-gate claim \\
Recorder & Clocked structured event writer & Connected, non-exclusive & Multiple writers remain; no single-log replay claim \\
\bottomrule
\end{tabular}
\caption{Implementation state in the audited repository.  K1--K5 completion
does not imply that every subsystem or output format is isolated.}
\label{tab:migration-state}
\end{table}

\subsection{Failure and reproducibility boundaries}

A cognition-stage exception propagates because the stage is control flow.
Bus-handler and plugin-setup failures instead warn and normally leave the
remaining run active; missing dependencies deactivate the affected plugin.
Operational resilience therefore makes a completed process different from a
complete treatment.  A reproducible artifact must include resolved stage
order, active plugin identifiers, warnings, configuration, source snapshot,
model/provider conditions, and specialized outputs.

The scoped K1--K5 migration is complete, but shared dictionaries, the context
escape hatch, direct domain calls, and multiple recording paths remain.
``K1--K5 complete'' describes completion of the migration plan, not universal
subsystem isolation or validation.
